% ============================================================================
% PART IV: THE THEORY IN FULL (author-written; precision-critical)
% ============================================================================
\part{The Theory in Full}

This part assembles the foundations of Part~I into the four results that
organize the \nm{} account: the episode-transfer lemma controlling the
invariance gap, the eviction-rate lemma controlling the forgetting
deficit, the decomposition proposition that yields the $2{\times}2$
prediction (including its super-additive signature), and the break-even
proposition that prices retention against carrying costs. We close with
the Gaussian construction that exhibits ERM's across-episode failure in
closed form, and with the honest remarks on what the coupled competition
dynamics are \emph{not} proven to do. Throughout, the register is that of
the GAUSE companion theory \cite{li2026gause}: idealized-model results
that explain mechanisms and generate falsifiable orderings, with the
quantitative weight carried by the controlled experiments of Part~V.

\section{Assumptions, Stated and Audited}
\label{sec:dd4-assumptions}

\begin{definition}[Episode risk]
\label{def:dd4-risk}
For regime $r$, episode environment $\xi$ with law $P_\xi$ on
$\mathcal{X} \times \mathcal{Y}$, and head $w$, the episode risk is
$L_\xi(w) = \E_{(x,y) \sim P_\xi}[\rho(Xw;\, y)]$, with $\rho$ the
decision regret of Part~I, $\rho \in [0, B]$,
$B = 2 k\, w_{\max}\, y_{\max}$.
\end{definition}

\begin{definition}[Assumption A1: episode exchangeability]
\label{def:dd4-a1}
For each regime $r$ there is a hyper-distribution $Q_r$ over episode
environments such that the historical episodes
$\xi_{r,1}, \dots, \xi_{r,E_r}$ and the deployed episode $\xi_{r,e'}$
are i.i.d.\ draws from $Q_r$. Write
$\mu_r(w) = \E_{\xi \sim Q_r} L_\xi(w)$ and
$\sigma_r^2(w) = \Var_{\xi \sim Q_r} L_\xi(w)$.
\end{definition}

\begin{warning}
A1 idealizes twice. \emph{Independence} fails when macro memory links
episodes (the 2011 crisis is causally downstream of 2008); \emph{identical
distribution} fails under secular drift (market microstructure in 1998
differs from 2020 in ways no regime label captures). We use A1 exactly as
Inv-PnCO uses its environment assumption \cite{yan2026invpnco}: as the
formal rendering of ``the next crisis is a fresh draw from the crisis
distribution,'' with the experiments designed so that A1 holds by
construction (the synthetic generator) and the limitations section of the
companion paper carrying the real-world caveat. Nothing in this part
silently strengthens A1.
\end{warning}

\begin{definition}[Assumption A2: dormancy and foreign updates]
\label{def:dd4-a2}
While regime $r$ is dormant for $D$ steps, an allocation mechanism $A$
routes learning. $A$ is \emph{reward-driven} if the specialist holding
$r$ is made to learn a currently active regime with per-step probability
at least $p > 0$; $A$ is \emph{reward-independent} if the assignment of
$r$ is measurable with respect to information independent of the
realized reward stream and, at convergence, routes to $r$'s owner only
regimes in its fixed niche set.
\end{definition}

The $p > 0$ clause is not vacuous bookkeeping: it is verified \emph{by
construction} for each reward-driven arm in our inventory (the monolith
learns the active regime with probability $1$; the router's
$\epsilon$-greedy exploration plus reward-chasing gate gives
$p \ge \epsilon / N$; the recent-performance allocator's softmax floor
gives $p \ge v_{\min} > 0$), and it is exactly the quantity a fund can
audit in its own allocator: \emph{what fraction of a losing desk's
capacity is reassigned per period?}

\section{The Invariance Gap}
\label{sec:dd4-invgap}

\begin{lemma}[Episode transfer]
\label{lem:dd4-transfer}
Under A1, for any fixed head $w$ and any $\delta \in (0, 1]$, with
probability at least $1 - \delta$ over $\xi_{e'} \sim Q_r$,
\[
  L_{\xi_{e'}}(w) \;\le\; \mu_r(w) +
  \sigma_r(w)\,\sqrt{\tfrac{1-\delta}{\delta}}
  \;\le\; \mu_r(w) + \frac{\sigma_r(w)}{\sqrt{\delta}} .
\]
\end{lemma}

\begin{proof}
Cantelli's inequality (Part~I) applied to the real random variable
$Z = L_\xi(w)$, $\xi \sim Q_r$, which has mean $\mu_r(w)$ and variance
$\sigma_r^2(w)$: for $\tau > 0$,
$\Pr[Z - \mu_r \ge \tau] \le \sigma_r^2 / (\sigma_r^2 + \tau^2)$.
Setting the right side to $\delta$ and solving,
$\tau = \sigma_r \sqrt{(1-\delta)/\delta}$, which is at most
$\sigma_r / \sqrt{\delta}$. No boundedness or distributional shape is
used; only the existence of the second moment, guaranteed by
$Z \in [0, B]$.
\end{proof}

\begin{keypoint}
The bound is monotone in \emph{both} $\mu_r(w)$ and $\sigma_r(w)$, and in
nothing else. Any training objective that drives the pair
$(\mu_r, \sigma_r)$ toward the lower-left of its Pareto frontier
optimizes the guarantee; the \nm{} objective
$\Var_e + \beta\,\E_e$ is the Lagrangian scalarization of exactly that
frontier, with $\beta$ the exchange rate. For each tail level $\delta$
there is a $\beta(\delta)$ whose minimizer over any convex model class
attains the frontier point optimizing the Lemma's bound (standard
scalarization; for non-convex classes the scalarization reaches the
convex hull of the frontier, a caveat we inherit silently from every
penalty method in the invariance literature \cite{krueger2021rex}).
\end{keypoint}

\begin{lemma}[Estimation price of few episodes]
\label{lem:dd4-estimation}
Fix $w$. With $E_r$ training episodes drawn i.i.d.\ from $Q_r$ and
$\hat\mu_r(w)$, $\hat\sigma_r(w)$ the empirical mean and standard
deviation of $\{L_{\xi_e}(w)\}_{e=1}^{E_r}$: for any
$\delta \in (0, 1)$, with probability at least $1 - \delta$,
\[
  |\hat\mu_r(w) - \mu_r(w)| \;\le\; B \sqrt{\frac{\log(2/\delta)}{2 E_r}},
  \qquad
  |\hat\sigma_r(w) - \sigma_r(w)| \;\le\; B \sqrt{\frac{2 \log(2/\delta)}{E_r - 1}} .
\]
\end{lemma}

\begin{proof}
The mean bound is Hoeffding's inequality (Part~I) for i.i.d.\
$[0, B]$-valued variables. For the deviation bound, note that the
empirical standard deviation of $[0,B]$-bounded variables is a function
of the sample with bounded differences: replacing one $L_{\xi_e}$ changes
$\hat\sigma_r$ by at most $B / \sqrt{E_r - 1}$ (each summand
$(L_e - \hat\mu)^2$ moves by at most $B^2$ terms damped by the
$1/(E_r{-}1)$ normalization; a one-line calculation, or see
\cite{maurer2009empirical} for the sharper empirical-Bernstein form).
McDiarmid's inequality then gives the stated tail, and
$\E[\hat\sigma_r] \le \sigma_r$ plus the same concentration controls the
bias direction.
\end{proof}

\begin{warning}[title={Important: the plug-in caveat}]
Lemmas~\ref{lem:dd4-transfer}--\ref{lem:dd4-estimation} hold for
\emph{fixed} $w$. The trained head $\hat w_r$ is a function of the
training episodes, so applying them at $w = \hat w_r$ is a plug-in
heuristic, not a uniform-convergence guarantee; making it uniform costs
a complexity term over the head class (linear heads in $d = 20$
dimensions would admit a standard covering-number argument) that we do
not need for the mechanism claims and therefore do not carry. This is
the same register at which Inv-PnCO deploys its Theorem~1, and we flag
rather than hide it. The practical content survives untouched: with
$E_r = 6$ episodes and $B = 0.4$, the mean's estimation radius at
$\delta = 0.1$ is $B\sqrt{\log 20 / 12} \approx 0.20$---half the regret
range---which is the formal statement of \emph{why} crisis-grade
invariance cannot be certified from six episodes alone and why the
sub-episode and cross-asset enlargements of Part~II are design defaults
rather than afterthoughts.
\end{warning}

\begin{corollary}[Bridge to divergence-style guarantees]
\label{cor:dd4-bridge}
If a representation $\Phi$ renders episode laws $\varepsilon$-invariant
in KL---$\KL(P_\xi(\cdot \mid \Phi)\, \|\, \bar P_r(\cdot \mid \Phi))
\le \varepsilon$ for $Q_r$-almost-every $\xi$---then every head $w$
factoring through $\Phi$ satisfies
$\sup_\xi |L_\xi(w) - \bar L_r(w)| \le B\sqrt{\varepsilon/2}$, hence
$\sigma_r(w) \le B\sqrt{\varepsilon/2}$, and
Lemma~\ref{lem:dd4-transfer} applies with that variance.
\end{corollary}

\begin{proof}
Immediate from the divergence-to-regret bridge of Part~I
($|L_P - L_{P'}| \le B\,\TV \le B\sqrt{\KL/2}$, Pinsker) applied
between each $P_\xi$ and the mixture $\bar P_r$; a uniformly
$c$-close-to-mean random variable has standard deviation at most $c$.
\end{proof}

This corollary is the formal connective tissue between the Inv-PnCO
guarantee language (divergences between solution distributions
\cite{yan2026invpnco}) and the regret currency of our decomposition: one
Pinsker step, at the explicit and admittedly loose price $B$.

\section{The Forgetting Deficit}
\label{sec:dd4-forget}

\begin{lemma}[Eviction rate under reward-driven allocation]
\label{lem:dd4-eviction}
Hard memory, capacity $K$; during $r$'s dormancy let $W_a \ge K$ distinct
regimes be active, let each dormant step independently deliver a foreign
learning assignment to $r$'s owner with probability $\ge p$, and
conditional on an assignment let the regime be drawn from the active set
with each regime's probability $\ge q_{\min} > 0$. Let $U_D$ be the
number of \emph{distinct} previously-unheld regimes learned in $D$ steps.
Then $r$'s head (and its episode buffer) is evicted on
$\{U_D \ge K\}$, and
\[
  \Pr[U_D < K] \;\le\;
  \Pr\!\big[\mathrm{Binom}(D,\ p\, q_{\min} (W_a - K + 1)) < K\big]
  \;\le\; \exp\!\Big(\!-2 D \big(p\,q_{\min}(W_a{-}K{+}1) -
  \tfrac{K}{D}\big)_+^2\Big),
\]
so $\Pr[\text{eviction}] \to 1$ exponentially in $D$. Moreover the
expected eviction time is bounded by the coupon-collector sum
\[
  \E[\text{time to } U = K] \;\le\; \frac{1}{p}
  \sum_{j=0}^{K-1} \frac{1}{q_{\min}\,(W_a - j)} .
\]
Under the soft model with per-foreign-update decay $\rho_{\mathrm{int}}$,
the retained strength after $N_D$ foreign updates is
$(1 - \rho_{\mathrm{int}})^{N_D}$, and $N_D \ge \mathrm{Binom}(D, p)$
stochastically, so $\E[\text{strength}] \le
(1 - p\,\rho_{\mathrm{int}})^{D} \to 0$.
\end{lemma}

\begin{proof}
While fewer than $K$ distinct foreign regimes have been acquired, at
least $W_a - K + 1$ of the active regimes are still unheld, so each step
delivers a \emph{new} distinct foreign regime with probability at least
$p\, q_{\min} (W_a - K + 1)$. The count of such successes in $D$ steps
thus stochastically dominates the stated binomial while $U < K$, giving
the first inequality; the second is Hoeffding's bound for the binomial
lower tail. The waiting-time bound is the sum of the $K$ geometric means
with success probabilities $p\, q_{\min}(W_a - j)$, $j = 0, \dots,
K{-}1$ (the coupon-collector computation of Part~II). LRU evicts $r$
precisely when $K$ distinct foreign heads have displaced it, since the
dormant $r$ is by definition the least recently used. For the soft
model, each foreign update multiplies $r$'s strength by
$(1 - \rho_{\mathrm{int}})$; taking expectations over the binomial count
gives the product form via $\E[(1-\rho)^{\mathrm{Binom}(D,p)}] =
(1 - p\rho)^D$.
\end{proof}

\begin{example}[Numbers at our calibration]
$K = 2$, $W_a = 3$ active common regimes, $p \approx 1$ (the monolith
always learns the active regime), $q_{\min} \approx 1/3$ at block scale:
the expected eviction time is on the order of two block lengths
($\approx 50$--$100$ days), and at crisis dormancy $D \approx 280$ days
the eviction probability is numerically indistinguishable from $1$. This
is why the E3 dormancy sweep finds the retention gap already saturated
at $D = 63$--$126$: the eviction event is binary, happens early in the
dormancy, and cannot happen twice.
\end{example}

\begin{proposition}[Dichotomy: reward-independence zeroes the deficit]
\label{prop:dd4-dichotomy}
Define the forgetting deficit
$\Phi(D; A) = \Pr[\text{$r$'s head lost at reactivation under } A]$
(hard) or $1 - \E[\text{retained strength}]$ (soft). Then:
(i) for any reward-driven $A$ satisfying A2,
$\Phi(D; A) \to 1$ as $D \to \infty$ at the rate of
Lemma~\ref{lem:dd4-eviction}; (ii) for any reward-independent covering
$A$, $\Phi(D; A) = 0$ for all $D$ in the hard model, and in the soft
model the owner's strength in $r$ decays only through the regime's own
drift, with \emph{zero} contribution from foreign updates.
\end{proposition}

\begin{proof}
(i) is the lemma. (ii): coverage gives an owner holding $r$ within
capacity; reward-independence means the assignment map does not respond
to the realized reward stream, so during $r$'s dormancy the owner is
assigned only regimes in its fixed niche set. If its niche set is
$\{r\}$ (the converged one-per-regime case) it receives no updates at
all; if its set contains other regimes, those are by construction
within capacity, so no eviction is triggered ($U_D = 0$ counts only
\emph{unheld} acquisitions) and---in the implemented soft model, where
decay attaches to foreign-\emph{regime} updates---no interference is
triggered either. The head and the episode buffer at reactivation are
bit-identical to their state at dormancy onset.
\end{proof}

\begin{keypoint}
The dichotomy is qualitative, not quantitative: the two classes are
different \emph{functions} of $D$ (one $\to 1$, one $\equiv 0$), which
is what licenses the experimental design---arm classes, not parameter
sweeps, separate them, and the E1 table confirms the class boundary at
$p \sim 10^{-6}$ while E3 confirms the saturation shape. Coverage is the
second conjunct and not decorative: random fixed niches are
reward-independent yet retain only what they own, and their doubled seed
variance in E1 is the measured price of leaving coverage to chance.
\end{keypoint}

\section{The Decomposition and the $2{\times}2$}
\label{sec:dd4-decomp}

\begin{definition}[Relearning cost]
\label{def:dd4-crelearn}
$C_{\mathrm{relearn}}$ is the expected excess of a from-prior head's
mean probe-window regret over a converged retained head's, for the same
regime and probe length $N_p$. It is an empirical constant of the
learning problem (model class, optimizer, signal-to-noise), not a
theorem; the E6 audit shows it collapsing as signal strength rises,
which is the honest reason the entire mechanism has a signal-to-noise
boundary.
\end{definition}

\begin{proposition}[Post-reactivation regret decomposition]
\label{prop:dd4-decomp}
Under A1--A2, when regime $r$ reactivates after dormancy $D$ into a
fresh episode, served under allocation $A$ by a head trained on $E_r$
episodes, then at tail level $\delta$, with the convention that the
invariance terms are evaluated at the served retained head $\hat w_r$,
\[
  \E[\rho_{\mathrm{react}}(r)] \;\le\;
  \underbrace{\hat\mu_r(\hat w_r) + \frac{\hat\sigma_r(\hat w_r)}{\sqrt\delta}
  + O\!\Big(B\sqrt{\tfrac{\log(1/\delta)}{E_r}}\Big)}_{\Ginv}
  \;+\;
  \underbrace{\Phi(D; A)\; C_{\mathrm{relearn}}}_{\Gforget},
\]
where the $O(\cdot)$ collects the estimation radii of
Lemma~\ref{lem:dd4-estimation}.
\end{proposition}

\begin{proof}
Condition on the retention event $R$ ($\Pr[R] = 1 - \Phi$). On $R$ the
served head is $\hat w_r$; Lemma~\ref{lem:dd4-transfer} bounds its
fresh-episode risk by $\mu_r + \sigma_r/\sqrt\delta$, and
Lemma~\ref{lem:dd4-estimation} converts population to empirical
quantities at the stated radii. On $R^c$ ($\Pr = \Phi$) the served head
is fresh; by Definition~\ref{def:dd4-crelearn} its probe regret exceeds
the retained benchmark by $C_{\mathrm{relearn}}$ in expectation, and the
benchmark itself is bounded by the same invariance terms. Total
expectation: $\E[\rho_{\mathrm{react}}] \le (1-\Phi)\,\Ginv + \Phi\,
(\Ginv + C_{\mathrm{relearn}}) = \Ginv + \Phi\, C_{\mathrm{relearn}}$.
\end{proof}

\begin{proposition}[Independent controls; super-additivity where it
matters]
\label{prop:dd4-additive}
$\Gforget$ is a function of $(A, D, K, \text{memory model})$ alone and
is zeroed by any reward-independent covering $A$
(Proposition~\ref{prop:dd4-dichotomy}) regardless of the training
objective; $\Ginv$ is a function of $(\text{objective}, Q_r, E_r)$ alone
and is optimized by the variance-penalized objective
(Lemma~\ref{lem:dd4-transfer} discussion) regardless of $A$. The
\emph{bound} is therefore additive in the two design axes. The
\emph{realized} interaction is predicted super-additive in one specific
cell: under reward-driven $A$ with $\Phi \approx 1$, the served head at
reactivation is fresh whatever objective trained the evicted one, so
improving $\Ginv$ moves nothing observable---the invariance axis is
inert in the reward-driven row.
\end{proposition}

\begin{remark}[The measured interaction]
Define $I = (\text{A1} - \text{A5}) + (\text{A1} - \text{A7}) -
(\text{A1} - \text{A6})$ on post-reactivation regret: $I = 0$ is
additivity, $I < 0$ is synergy. The companion experiments measure
$I = -0.0015 \pm 0.0005$ per-seed (95\% CI excluding zero), realized
almost entirely through $\text{A1} - \text{A7} \approx 0$
($p = 0.88$)---the inert cell, exactly where
Proposition~\ref{prop:dd4-additive} puts it. We are careful about
logical direction: the additive \emph{bound} motivates but does not
entail the super-additive \emph{realization}; the prediction that A7
specifically would sit at A1's level was pre-registered from the
mechanism (an evicted head serves nobody) and is the account's sharpest
falsifiable signature---a ``both tricks help'' story without the
two-clock structure would have A7 strictly between A1 and A6.
\end{remark}

\section{Break-Even Dormancy: Pricing Retention}
\label{sec:dd4-breakeven}

\begin{proposition}[The retention window]
\label{prop:dd4-breakeven}
Charge carrying cost $c$ per idle specialist-day and let $V$ convert one
unit of mean probe regret to value over the $N_p$-day probe. Over one
dormancy--reactivation cycle of length $D$, retaining $r$'s specialist
pays iff
\[
  c\, D \;\le\; \Phi(D; A_{\mathrm{rd}})\;
  C_{\mathrm{relearn}}\, N_p\, V .
\]
Since $\Phi$ is increasing and saturates at $1$ while $cD$ grows
linearly, the solution set is an interval $[D_{\min}, D_{\max}]$
(possibly empty): $D_{\min}$ where $\Phi$ ramps (eviction becomes likely
enough to fear), $D_{\max} = C_{\mathrm{relearn}} N_p V / c$ where
cumulative carrying cost overtakes the one-shot saving. The window is
empty iff $c > \max_D \Phi(D) C_{\mathrm{relearn}} N_p V / D$, i.e.,
for carrying costs above a threshold set entirely by measurable
quantities.
\end{proposition}

\begin{example}[At the E3 calibration]
$C_{\mathrm{relearn}} \approx 0.004$ regret per probe day, $N_p = 15$:
the cycle saving is $\approx 0.06\, V$. At a carrying cost of $1\%$ of
the daily probe saving ($c = 4 \times 10^{-5} V$), $D_{\max} \approx
1{,}500$ trading days and $D_{\min} \approx 40$--$60$ (where E3 shows
$\Phi$ ramping); at $10\%$, $D_{\max} \approx 150$ days---tighter than
one crisis cycle, so a fund at that cost level should \emph{not} retain
a dedicated crisis specialist and should plan to relearn; the indifference
threshold sits near $15\%$. The proposition's value is not a universal
verdict but the conversion of an organizational argument
(``keep the crisis desk?'') into three measurable numbers.
\end{example}

\section{What the Competition Dynamics Are Not Proven to Do}
\label{sec:dd4-honesty}

Two honest closings. First, \emph{we prove no convergence theorem for
the coupled winner-take-all dynamics}: specialists' window wins depend on
heads, heads on training, training on wins---an endogenous feedback for
which adversarial-regret guarantees of the EG update do not apply (Part~I)
and for which we know of no clean fixed-point argument covering the full
mechanism. What we have is the GAUSE precedent (the same dynamics
converge to covering one-per-regime assignments across all its domains
and protocols \cite{li2026gause}), our own twenty-for-twenty convergence
record with the crisis-regime affinity wrinkle reported, and the
single-window concentration remark---whose pre-registered implication
(batching necessary) the $W_c$ ablation \emph{refuted}, teaching the
corrected lesson that cross-window EG accumulation is the operative
averager. Second, \emph{the decomposition's two terms are bounds, not
identities}: a mechanism could in principle beat the bound's additivity
(and ours measurably does, in the predicted direction and cell). The
theory's role is to say which knob owns which term and where each term
vanishes---claims the eleven-arm design tests directly---not to predict
regret to the fourth decimal. Where the theory and the experiments
disagreed, this document reports the experiment and revises the story,
twice (batching; the high-SNR prototype that became E6).
